% !TeX program = xelatex
\documentclass[12pt,a4paper]{article}
\usepackage[margin=2cm]{geometry}
\usepackage{fontspec}
\usepackage{amsmath,amssymb,enumitem}
\setmainfont{Times New Roman}
\setlength{\parindent}{0pt}\setlength{\parskip}{6pt}
\begin{document}
\begin{center}{\Large\bfseries Đáp án Toán vào lớp 10 Hà Nội năm 2014\\Phần Đại số}\end{center}
\textbf{Câu 26.} Thay $x=9$:
\[
A=\frac{\sqrt9+1}{\sqrt9-1}=\frac{3+1}{3-1}=2.
\]

\textbf{Câu 27.}
\begin{enumerate}[label=\alph*)]
\item Với $x>0$, $x\ne1$,
\begin{align*}
P&=\left(\frac{x-2}{\sqrt{x}(\sqrt{x}+2)}+\frac{1}{\sqrt{x}+2}\right)\cdot\frac{\sqrt{x}+1}{\sqrt{x}-1}\\
&=\frac{x-2+\sqrt{x}}{\sqrt{x}(\sqrt{x}+2)}\cdot\frac{\sqrt{x}+1}{\sqrt{x}-1}\\
&=\frac{(\sqrt{x}-1)(\sqrt{x}+2)}{\sqrt{x}(\sqrt{x}+2)}\cdot\frac{\sqrt{x}+1}{\sqrt{x}-1}
=\frac{\sqrt{x}+1}{\sqrt{x}}.
\end{align*}
\item Put $t=\sqrt{x}>0$ (and $t\ne1$). Then
\[
2\frac{t+1}{t}=2t+5
\iff 2t+2=2t^2+5t
\iff 2t^2+3t-2=0.
\]
Thus $t=\frac12$ or $t=-2$. Since $t>0$, $t=\frac12$, hence $x=\frac14$.
\end{enumerate}

\textbf{Câu 28.} Gọi $x$ là số sản phẩm mỗi ngày theo kế hoạch, với $x\in\mathbb N$ và $0<x<1100$. Số ngày theo kế hoạch là $1100/x$, còn khi tăng sản lượng mỗi ngày thêm $5$ thì số ngày là $1100/(x+5)$. Vì hoàn thành sớm $2$ ngày:
\[
\frac{1100}{x+5}+2=\frac{1100}{x}
\iff x^2+5x-2750=0
\iff x=50\ \text{hoặc}\ x=-55.
\]
Loại $x=-55$ vì không phù hợp ý nghĩa; vậy mỗi ngày theo kế hoạch sản xuất $50$ sản phẩm.

\textbf{Câu 29.} Điều kiện $x+y\ne0$, $y\ne1$. Đặt $X=1/(x+y)$ và $Y=1/(y-1)$. Khi đó
\[
\begin{cases}4X+Y=5,\\X-2Y=-1\end{cases}
\quad\Rightarrow\quad X=1,\ Y=1.
\]
Suy ra $x+y=1$ và $y-1=1$, nên $y=2$, $x=-1$. Nghiệm là $(-1;2)$.

\textbf{Câu 30.}
\begin{enumerate}[label=\alph*)]
\item Hoành độ giao điểm thỏa mãn $x^2=-x+6$, tức $(x+3)(x-2)=0$. Với $x=2$ thì $y=4$; với $x=-3$ thì $y=9$. Do đó hai giao điểm là $A(2;4)$ và $B(-3;9)$.
\item Diện tích tam giác với hai vectơ $\overrightarrow{OA}=(2,4)$, $\overrightarrow{OB}=(-3,9)$:
\[
S_{OAB}=\frac12\left|2\cdot9-4\cdot(-3)\right|=15.
\]
\end{enumerate}

\textbf{Câu 32.} Do $a+b+c=2$,
\[
2a+bc=a(a+b+c)+bc=(a+b)(a+c).
\]
Theo bất đẳng thức AM--GM,
\[
\sqrt{2a+bc}=\sqrt{(a+b)(a+c)}\le\frac{(a+b)+(a+c)}2.
\]
Tương tự cho hai căn còn lại. Cộng ba bất đẳng thức được
\[
Q\le\frac{(a+b+a+c)+(b+c+b+a)+(c+a+c+b)}2
=2(a+b+c)=4.
\]
Dấu bằng xảy ra khi $a=b=c=2/3$, thỏa mãn điều kiện. Vậy $\max Q=4$.
\end{document}
